\documentclass[book.tex]{subfiles}

\begin{document}
\chapter{Quantum Fourier Transform}
\label{chap:quantum_fourier_transform}
\noindent\rule{11cm}{0.4pt} \\
\textbf{Prerequisites:} \autoref{chap:quantum_circuits} \\
\textbf{Difficulty Level:} ** \\
\noindent\rule{11cm}{0.4pt}
\vspace{5mm}

The quantum Fourier transform implements the Fourier transform as a quantum circuit and is a fundamental part of many circuits. The quantum fourier transform acts on an input state
\begin{align}
    |X\rangle &= \sum_{j=0}^{N-1} x_j |j \rangle
\end{align}
and maps it to the state
\begin{align}
|Y\rangle &= \sum_{k=0}^{N-1} y_k |k\rangle
\end{align}
where we have that
\begin{align}
y_k &= \frac{1}{\sqrt{N}} \sum_{j=0}^{N-1} x_j \omega_N^{jk}
\end{align}

Intuitively, the QFT acts as a change of basis transformation mapping from the a standard computational basis to a fourier basis for the quantum system. 
\end{document}
